\documentclass[acmsmall,nonacm,screen]{acmart}

\usepackage{mathpartir}
\usepackage{booktabs}

\citestyle{acmauthoryear}
\settopmatter{printacmref=false}

\AtEndPreamble{\theoremstyle{acmdefinition}\newtheorem{remark}[theorem]{Remark}}

\newcommand{\kw}[1]{\mathsf{#1}}
\newcommand{\ret}{\kw{ret}}
\newcommand{\letin}[3]{\kw{let}\;#1 = #2\;\kw{in}\;#3}
\newcommand{\ite}[3]{\kw{if}\;#1\;\kw{then}\;#2\;\kw{else}\;#3}
\newcommand{\chk}[2]{\kw{check}_{#1}\;#2}
\newcommand{\vote}[2]{\kw{vote}_{#1}\;#2}
\newcommand{\orelse}{\mathbin{\kw{orelse}}}
\newcommand{\retry}[2]{\kw{retry}_{#1}\;#2}
\newcommand{\rty}[2]{\{#1 \mid #2\}}
\newcommand{\at}[2]{\mathrel{@}(#1, #2)}
\newcommand{\sem}[1]{[\![#1]\!]}
\newcommand{\Bad}{\mathrm{Bad}}
\newcommand{\Fail}{\mathrm{Fail}}
\newcommand{\cO}{\mathcal{O}}
\newcommand{\nm}[1]{\mathit{#1}}
\newcommand{\tl}[3]{T_{#1}^{#2}(#3)}

% written by grades.py; do not edit
\newcommand{\exCases}{55}
\newcommand{\exChainTen}{40.1}
\newcommand{\exGapBound}{5/12}
\newcommand{\exGapExact}{3/8}
\newcommand{\exLoopCalls}{6.5}
\newcommand{\exLoopTen}{34.9}
\newcommand{\exPatchCalls}{4.2}
\newcommand{\exPatchFive}{12.4}
\newcommand{\exPatchLimit}{16.7}
\newcommand{\exPatchOne}{4.0}
\newcommand{\exPrograms}{3,000}
\newcommand{\exRules}{12}
\newcommand{\exSlowest}{50}
\newcommand{\exTightBad}{68}
\newcommand{\exTightCost}{78}
\newcommand{\exTightFail}{73}
\newcommand{\exUnionBad}{60}
\newcommand{\exUnionFail}{63}
\newcommand{\exVoteBadEleven}{0.753}
\newcommand{\exVoteBadFive}{0.683}
\newcommand{\exVoteBadOne}{0.600}
\newcommand{\exVoteBadThree}{0.648}
\newcommand{\exVoteBadTwentyFive}{0.846}
\newcommand{\exVoteEleven}{0.078}
\newcommand{\exVoteFive}{0.163}
\newcommand{\exVoteOne}{0.300}
\newcommand{\exVoteThree}{0.216}
\newcommand{\exVoteTwentyFive}{0.017}


\begin{document}

\title{Turning Errors into Abstentions: Types for Programs That Generate and Check}

\author{Claude-Fable-5.2}
\authornote{An AI model. The proofs are on paper and have not been checked by a machine. A
script computes the exact worst case of small programs and tests the theorems against it. At
the time of writing nobody has reviewed the paper. This version replaces a shorter one; the
record of the paper lists the changes.}

\begin{abstract}
Programs now call oracles, such as language models, whose answers are right only with some
probability, and they wrap these calls in checkers, retries, and votes. We give a type system
for such programs. A type carries two bounds: $\varepsilon$ on the probability of returning a
wrong value, and $\gamma$ on the probability of not returning a right one. An oracle call
introduces error. A sound checker sets $\varepsilon$ to zero, so it turns errors into
abstentions, and a retry multiplies $\gamma$, so it turns abstentions into further calls. An imperfect checker, such as a test suite or a second model, multiplies
$\varepsilon$ by its rate of false acceptance. The system is sound against adaptive oracles:
the bounds hold for every environment whose failure probability is bounded at each call,
whatever it has seen before. The exact worst case of a small program is the value of a game
against such an environment. We compute it, and we use it to test the theorem, to show that
each rule is needed, and to measure how tight the grades are. They are exact for the standard
patterns and for about two thirds of random programs, and we show why no pair of numbers can
be exact in general. Three consequences stand out. A loop that generates until a sound check
passes never returns a wrong value. A loop that generates until an imperfect check passes
converges to the false discovery rate of the checker, which for a weak generator is far above
the checker's own error. A majority vote keeps its binomial guarantee against adaptive
oracles, but it needs unique answers, and it never reaches zero.
\end{abstract}

\keywords{probabilistic programs, refinement types, graded types, oracles, language models, checkers}

\authorsaddresses{}

\maketitle

\section{Introduction}

A common way to get reliable results from an unreliable source is to check them. A language
model proposes a proof and a kernel checks it. A model proposes a patch and a test suite runs.
A model extracts a date and a parser validates its format. When the check fails, the program
tries again or gives up. Where no check exists, the program may ask several times and take
the majority~\citep{wang2022}. We call such programs \emph{generate-and-check} programs.

Their guarantees are usually stated informally, and two different things are confused. A
\emph{wrong} result is a value that violates the specification. An \emph{abstention} is the
absence of a result. A checker does not make an oracle right more often. It converts wrong
results into abstentions, which the caller can see and handle. A retry does not reduce wrong
results either. It reduces abstentions, at the price of more calls.

This paper makes the distinction a matter of types. A computation has a type
$\rty{x}{\varphi} \at{\varepsilon}{\gamma}$: its result satisfies $\varphi$, except that with
probability at most $\varepsilon$ it is a wrong value, and with probability at most $\gamma$ it
is not a right value, which includes abstaining. The typing rules say how each construct moves
the two bounds. Table~\ref{tab:rules} lists them.

\begin{table}[t]
\caption{What each construct does to the bound $\varepsilon$ on wrong results and to the bound
$\gamma$ on failures. We write $a \oplus b$ for $a + b - ab$, and $\tl{k}{m}{p}$ for the
probability of at least $m$ successes in $k$ independent trials of rate $p$.}
\label{tab:rules}
\small
\begin{tabular}{@{}lll@{}}
\toprule
Construct & Wrong results & Failures \\
\midrule
oracle with failure rate $\delta$ & $\delta$ & $\delta$ \\
sequence $\kw{let}$ & $\varepsilon_1 \oplus \varepsilon_2$ & $\gamma_1 \oplus \gamma_2$ \\
exact checker & $0$ & $\gamma$ \\
checker with false acceptance $\alpha$, false rejection $\beta$ & $\alpha\,\varepsilon$ & $\gamma \oplus \beta$ \\
fallback $\kw{orelse}$ & $\varepsilon_1 + (\gamma_1 - \varepsilon_1)\,\varepsilon_2$ & $\varepsilon_1 + (\gamma_1 - \varepsilon_1)\,\gamma_2$ \\
majority of $k$ runs & $\tl{k}{\lfloor k/2 \rfloor + 1}{\varepsilon}$ & $\tl{k}{\lceil k/2 \rceil}{\gamma}$ if the answer is unique \\
\bottomrule
\end{tabular}
\end{table}

\paragraph{Adaptive oracles.} A model does not fail independently. Its second answer depends
on its first, on the checker's report, and on everything else in its context. We therefore
make no assumption of independence. An environment may choose the distribution of each answer
as it likes, knowing the whole history, as long as each single answer is wrong with
probability at most the stated rate. Theorem~\ref{thm:sound} holds for all such environments.
The assumption that remains is that each call has a fresh chance to succeed, and
Section~\ref{sec:fresh} shows that the guarantees of retries and votes fail without it.

\paragraph{Exact worst cases.} Against an adaptive environment, the worst case of a program
is the value of a game: the environment moves at each call, and it wants a wrong result, or a
failure, or many calls. For programs over a small set of values the game can be solved by
backward induction. The artifact of this paper does that, with exact rational numbers, and
the paper uses it in three ways. It tests the soundness theorem on \exPrograms{} random
programs. It refutes each of \exRules{} unsound variants of the rules by a random program. And
it measures how tight the grades are: the bound on wrong results equals the value of the game
for \exTightBad{} percent of the random programs, and the bound on failures for \exTightFail{}
percent. The rules are exact for the patterns that programs use, such as the loop with a
checker and the vote. They cannot be exact in general, and Section~\ref{sec:tight} gives a
program on which any system that abstracts a computation by two numbers must lose.

\paragraph{Consequences.} Section~\ref{sec:consequences} reads the rules as laws. The loop
that generates until a sound check passes has no error at all, an abstention probability
$\delta^n$, and at most $1/(1-\delta)$ expected calls. With an imperfect checker the same loop
behaves differently: every retry is another chance for a wrong answer to slip through, and the
error of the loop rises with the number of attempts toward
$\delta\alpha / (\delta\alpha + (1-\delta)(1-\beta))$, the probability that an accepted answer
is wrong. For a generator that is right one time in five and tests that pass one wrong patch
in twenty, a single attempt returns a wrong patch with probability \exPatchOne{} percent, five
attempts with \exPatchFive{} percent, and the limit is \exPatchLimit{} percent. A majority
vote over $k$ answers has the binomial tail as its error even when the oracle adapts, but only
when right answers are unique and the oracle is right more often than not.

\paragraph{Contributions.}
\begin{itemize}
\item A language with oracles, checkers with two-sided error, fallbacks, and votes, and a
  semantics in which the environment is an adaptive adversary with bounded failure rates
  (Section~\ref{sec:lang}).
\item A type system with two grades and its soundness theorem, with a bound on the expected
  number of calls (Sections~\ref{sec:types} and~\ref{sec:soundness}).
\item The exact worst case as the value of a game, a measurement of the tightness of the
  grades, and an example that shows why two numbers cannot be exact (Section~\ref{sec:tight}).
\item Five consequences for programs that generate and check (Section~\ref{sec:consequences}).
\item An artifact that solves the game exactly and tests every claim above against it
  (Section~\ref{sec:tests}).
\end{itemize}

The result that a checked computation has error zero is the quantitative form of an old
principle: a proof checker makes the prover irrelevant to trust. What stays is the
specification that the checker enforces, which a human has to read~\citep{tpp1}, and the calls
that nothing checks.

\paragraph{Status of the claims.} The proofs in this paper are written out and have not been
mechanized. The artifact tests the theorems on small programs with exact arithmetic. It does
not estimate the failure rate of any real model: the rates are inputs.

\section{Overview: Three Loops}
\label{sec:overview}

\paragraph{A sound check.} A model with failure rate $\delta$ proposes a proof of a statement
$s$, and a kernel checks it:
\[
\retry{n}{(\chk{\nm{kernel}}{\nm{prove}(s)})} .
\]
The call has the grade $(\delta, \delta)$: it returns a wrong proof with probability at most
$\delta$. The kernel accepts no wrong proof, so the checked call has the grade $(0, \delta)$:
it never returns a wrong proof, and it abstains with probability at most $\delta$. A fallback
runs its second branch only when the first abstains, so $n$ attempts have the grade
$(0, \delta^n)$. For a model that is right one time in ten, ten attempts abstain with
probability at most \exLoopTen{} percent, make at most \exLoopCalls{} calls on average, and
return a wrong proof with probability zero. The failure rate of the model affects how often the loop gives
up and how long it runs, and nothing else.

\paragraph{An imperfect check.} A model proposes a patch, and a test suite runs. A test suite
is not a kernel: it passes some wrong patches~\citep{smith2015}. Let it accept a wrong patch
with probability at most $\alpha$. One attempt returns a wrong patch with probability
$\delta\alpha$. The loop tries again whenever the tests fail, which happens with probability
$r = \delta(1 - \alpha)$, and each further attempt is a further chance for a wrong patch to
pass. The grades give the error $\delta\alpha\,(1 + r + \cdots + r^{n-1})$ for $n$ attempts.
With $\delta = 0.8$ and $\alpha = 0.05$ this is \exPatchOne{} percent for one attempt,
\exPatchFive{} percent for five, and it converges to \exPatchLimit{} percent. A loop that runs
until the tests pass is more than three times as wrong as its tests.

\paragraph{A vote.} Where nothing can be checked, a program asks $k$ times and returns the
answer that more than half of the runs gave. If each answer is wrong with probability at most
$0.3$, a wrong answer wins the vote with probability at most \exVoteThree{} for $k = 3$,
\exVoteFive{} for $k = 5$, and \exVoteEleven{} for $k = 11$. This holds although the oracle
may coordinate its answers, because it must give each one a fresh chance. The vote differs
from the check in three ways: its error is never zero, it needs the right answer to be unique,
and it helps only an oracle that is right more often than wrong. At a failure rate of $0.6$
the same numbers are \exVoteBadThree, \exVoteBadFive{} and \exVoteBadEleven.

\section{Language}
\label{sec:lang}

Values $v$ are data. An outcome $r$ is a value or the abstention $\bot$. Expressions are
\[
e ::= \ret\;v \mid \letin{x}{e}{e} \mid \ite{b}{e}{e} \mid f(v) \mid o(v) \mid \chk{j}{e}
  \mid e \orelse e \mid \vote{k}{e} .
\]
Here $b$ is a Boolean test on values, $f$ ranges over total deterministic functions, $o$ over
oracle names, $j$ over checker names, and $k$ over positive numbers. The expression
$\chk{j}{e}$ runs $e$, shows its value to the checker $j$, and abstains unless $j$ accepts.
The expression $e_1 \orelse e_2$ runs $e_1$ and runs $e_2$ if $e_1$ abstains. The expression
$\vote{k}{e}$ runs $e$ for $k$ times and returns the value that more than half of the runs
returned; it abstains if there is none. Abstention propagates through $\kw{let}$. We write
$\retry{n}{e}$ for $e \orelse \cdots \orelse e$ with $n$ copies of $e$. The language has no
recursion, so every program makes a bounded number of calls.

\paragraph{Environments.}
A \emph{history} $h$ is the sequence of calls made so far, with their inputs and outputs. An
\emph{environment} $E$ assigns to each oracle $o$, input $v$, and history $h$ a distribution
$E(o, v, h)$ over values with finite support, and to each checker $j$, value $u$, and history
$h$ a probability $E(j, u, h)$ of acceptance. The answer of an oracle and the verdict of a
checker may thus depend on everything that happened before.

Each oracle $o$ comes with a specification: a relation $\psi_o(v, u)$ between inputs and right
outputs, and a failure rate $\delta_o \in [0, 1]$. Each checker $j$ comes with a predicate
$p_j$ on values and two rates: $\alpha_j$ for false acceptance and $\beta_j$ for false
rejection. The class $\cO$ consists of the environments that satisfy, for every $o$, $j$, $v$,
$u$, and $h$,
\[
\Pr_{u' \sim E(o, v, h)}[\neg \psi_o(v, u')] \le \delta_o, \qquad
E(j, u, h) \le \alpha_j \text{ if } \neg p_j(u), \qquad
1 - E(j, u, h) \le \beta_j \text{ if } p_j(u) .
\]
The bounds are required at each call and for each history. They are not an assumption of
independence. An environment in $\cO$ may choose its distributions adversarially.

A checker with $\alpha_j = \beta_j = 0$ is \emph{exact}: it is the deterministic predicate
$p_j$. A proof kernel and a parser are exact checkers. A checker with $\alpha_j = 0$ is
\emph{sound}. A test suite, a linter, and a second model that judges an answer are checkers
with positive rates.

\paragraph{Semantics.}
For a closed expression $e$, an environment $E$, and a history $h$, the semantics
$\sem{e}_E(h)$ is a distribution over pairs of an outcome and a history. It is defined by
recursion on $e$ in the usual way~\citep{kozen1981}. The expression $\ret\;v$ yields $(v, h)$.
The call $o(v)$ draws $u$ from $E(o, v, h)$ and yields $u$ with the history extended by this
call. The expression $\letin{x}{e_1}{e_2}$ draws $(r_1, h_1)$ from $\sem{e_1}_E(h)$ and yields
$(\bot, h_1)$ if $r_1 = \bot$, and otherwise continues with $\sem{e_2[r_1/x]}_E(h_1)$. The
expression $\chk{j}{e}$ draws $(r, h_1)$ from $\sem{e}_E(h)$. If $r = \bot$ it yields
$(\bot, h_1)$. Otherwise it yields $r$ with probability $E(j, r, h_1)$ and $\bot$ with the
remaining probability, in both cases with the history extended by the verdict. The expression
$e_1 \orelse e_2$ draws $(r_1, h_1)$ from $\sem{e_1}_E(h)$, yields it if $r_1 \neq \bot$, and
otherwise continues with $\sem{e_2}_E(h_1)$. The expression $\vote{k}{e}$ runs $e$ for $k$
times, each run from the history that the previous run left, and yields the value that occurs
more than $k/2$ times among the outcomes, or $\bot$.

\section{Types}
\label{sec:types}

A type $\rty{x}{\varphi}$ is a predicate on values. A context $\Phi$ is a set of predicates on
the free variables. We write $\Phi \models \varphi$ when every valuation that satisfies $\Phi$
satisfies $\varphi$. A \emph{grade} is a pair $(\varepsilon, \gamma)$ of numbers with
$0 \le \varepsilon \le \gamma \le 1$. Two operations on probabilities occur in the rules:
\[
a \oplus b = a + b - ab, \qquad
\tl{k}{m}{p} = \sum_{i=m}^{k} \binom{k}{i}\, p^i (1-p)^{k-i} .
\]
The first is the probability that at least one of two independent events of probabilities $a$
and $b$ occurs, and the second is the probability that at least $m$ of $k$ independent events
of probability $p$ occur. Both are monotone in their probabilities. The judgment
\[
\Phi \vdash e : \rty{x}{\varphi} \at{\varepsilon}{\gamma}
\]
is defined in Figure~\ref{fig:rules}. Each pure function $f$ has a specification $\varphi_f$
with $\varphi_f(v, f(v))$ for all $v$.

\begin{figure}[t]
\begin{mathpar}
\inferrule[Ret]{\Phi \models \varphi[v/x]}{\Phi \vdash \ret\;v : \rty{x}{\varphi} \at{0}{0}}
\and
\inferrule[Fun]{ }{\Phi \vdash f(v) : \rty{x}{\varphi_f(v, x)} \at{0}{0}}
\and
\inferrule[Oracle]{ }{\Phi \vdash o(v) : \rty{x}{\psi_o(v, x)} \at{\delta_o}{\delta_o}}
\and
\inferrule[Let]
  {\Phi \vdash e_1 : \rty{x}{\varphi_1} \at{\varepsilon_1}{\gamma_1} \\
   \Phi, \varphi_1 \vdash e_2 : \rty{y}{\varphi_2} \at{\varepsilon_2}{\gamma_2} \\
   x \notin \mathrm{fv}(\varphi_2)}
  {\Phi \vdash \letin{x}{e_1}{e_2} : \rty{y}{\varphi_2}
    \at{\varepsilon_1 \oplus \varepsilon_2}{\gamma_1 \oplus \gamma_2}}
\and
\inferrule[Check]
  {\Phi \vdash e : \rty{x}{\chi} \at{\varepsilon}{\gamma} \\
   \Phi, p_j(x) \models \varphi \\ \Phi, \chi \models p_j(x)}
  {\Phi \vdash \chk{j}{e} : \rty{x}{\varphi} \at{\alpha_j\,\varepsilon}{\gamma \oplus \beta_j}}
\and
\inferrule[Check-Any]
  {\Phi \vdash e : \rty{x}{\chi} \at{\varepsilon}{\gamma} \\ \Phi, p_j(x) \models \varphi}
  {\Phi \vdash \chk{j}{e} : \rty{x}{\varphi} \at{\alpha_j}{1}}
\and
\inferrule[OrElse]
  {\Phi \vdash e_1 : \rty{x}{\varphi} \at{\varepsilon_1}{\gamma_1} \\
   \Phi \vdash e_2 : \rty{x}{\varphi} \at{\varepsilon_2}{\gamma_2}}
  {\Phi \vdash e_1 \orelse e_2 : \rty{x}{\varphi}
    \at{\varepsilon_1 + (\gamma_1 - \varepsilon_1)\,\varepsilon_2}
       {\varepsilon_1 + (\gamma_1 - \varepsilon_1)\,\gamma_2}}
\and
\inferrule[Vote]
  {\Phi \vdash e : \rty{x}{\varphi} \at{\varepsilon}{\gamma} \\
   \Phi, \varphi, \varphi[x'/x] \models x = x'}
  {\Phi \vdash \vote{k}{e} : \rty{x}{\varphi}
    \at{\tl{k}{\lfloor k/2 \rfloor + 1}{\varepsilon}}{\tl{k}{\lceil k/2 \rceil}{\gamma}}}
\and
\inferrule[Vote-Any]
  {\Phi \vdash e : \rty{x}{\varphi} \at{\varepsilon}{\gamma}}
  {\Phi \vdash \vote{k}{e} : \rty{x}{\varphi} \at{\tl{k}{\lfloor k/2 \rfloor + 1}{\varepsilon}}{1}}
\and
\inferrule[If]
  {\Phi, b \vdash e_1 : \rty{x}{\varphi} \at{\varepsilon}{\gamma} \\
   \Phi, \neg b \vdash e_2 : \rty{x}{\varphi} \at{\varepsilon}{\gamma}}
  {\Phi \vdash \ite{b}{e_1}{e_2} : \rty{x}{\varphi} \at{\varepsilon}{\gamma}}
\and
\inferrule[Sub]
  {\Phi \vdash e : \rty{x}{\varphi} \at{\varepsilon}{\gamma} \\
   \Phi, \varphi \models \varphi' \\ \varepsilon \le \varepsilon' \\ \gamma \le \gamma'
   \\ \varepsilon' \le \gamma'}
  {\Phi \vdash e : \rty{x}{\varphi'} \at{\varepsilon'}{\gamma'}}
\end{mathpar}
\Description{Eleven typing rules with two grades.}
\caption{Typing rules.}
\label{fig:rules}
\end{figure}

\paragraph{Sequences.} In \textsc{Let} the grades compose with $\oplus$. The union bound
$\varepsilon_1 + \varepsilon_2$ would be sound as well, and the first version of this paper
used it. The product form is the exact worst case of two calls in a row, and it keeps grades
below $1$.

\paragraph{Checks.} The rule \textsc{Check} has two side conditions. The first says that what
the checker's predicate accepts satisfies $\varphi$. With it, a wrong result of the checked
expression is a wrong result of $e$ that the checker accepted, which has probability at most
$\alpha_j \varepsilon$; for a sound checker this is $0$. The second says that the predicate
holds of the values that $e$ was supposed to produce, so that a right value is rejected with
probability at most $\beta_j$. Without the second condition a right value may fall outside
the predicate, and \textsc{Check-Any} gives what is left: every accepted wrong value was
accepted against the odds $\alpha_j$, and the failure bound is trivial. A checker that
rejects everything is sound.

\paragraph{Fallbacks.} In \textsc{OrElse} the second computation runs only when the first
abstains. The first is wrong with probability at most $\varepsilon_1$ and fails with
probability at most $\gamma_1$, so it abstains with probability at most
$\gamma_1 - \varepsilon_1$ when it is wrong as often as it may be. A wrong value of the first
computation is returned as it is, which is why $\varepsilon_1$ appears in both bounds.

\paragraph{Votes.} A wrong value wins a vote only if more than half of the runs returned
wrong values, and \textsc{Vote} bounds the probability of that by a binomial tail. For the
failure bound the type must determine the value: if two different values can both be right,
right runs may disagree, and no value need reach a majority. \textsc{Vote-Any} is the rule
for that case.

\begin{lemma}
\label{lem:grades}
In every derivable judgment, $\varepsilon \le \gamma$.
\end{lemma}

\begin{proof}
By induction on the derivation. The operations $\oplus$ and $\tl{k}{m}{\cdot}$ are monotone,
$\tl{k}{m}{p}$ decreases in $m$, and $\lfloor k/2 \rfloor + 1 \ge \lceil k/2 \rceil$. For
\textsc{OrElse}, $\varepsilon_2 \le \gamma_2$ and $\gamma_1 - \varepsilon_1 \ge 0$.
\end{proof}

\section{Soundness}
\label{sec:soundness}

\subsection{The Theorem}

For a closed expression $e$ and a closed type $\rty{x}{\varphi}$, an outcome is \emph{bad}
when it is a value that violates $\varphi$, and it is a \emph{failure} when it is bad or
$\bot$. The judgment $\Phi \vdash e : \rty{x}{\varphi} \at{\varepsilon}{\gamma}$ is
\emph{valid} when, for every substitution $\theta$ of values for the free variables that
satisfies $\Phi$, every $E \in \cO$, and every history $h$,
\[
\Pr_{\sem{e\theta}_E(h)}[\Bad] \le \varepsilon
\qquad\text{and}\qquad
\Pr_{\sem{e\theta}_E(h)}[\Fail] \le \gamma .
\]

The rule for votes needs a fact about sequences of events that are not independent.

\begin{lemma}[Domination]
\label{lem:dom}
Let $X_1, \ldots, X_k$ be random variables with values in $\{0, 1\}$ such that, for every $i$
and every outcome of $X_1, \ldots, X_{i-1}$ and of anything else that is known before the
$i$-th variable is drawn, the conditional probability of $X_i = 1$ is at most $p$. Then
$\Pr[X_1 + \cdots + X_k \ge m] \le \tl{k}{m}{p}$ for every $m$.
\end{lemma}

\begin{proof}
By induction on $k$. For $k = 0$ both sides are $1$ if $m \le 0$ and $0$ otherwise. For
$k > 0$ let $q = \Pr[X_1 = 1] \le p$. After $X_1$ is drawn, the remaining variables satisfy
the hypothesis again, so by the induction hypothesis the probability of the event is at most
$q\,\tl{k-1}{m-1}{p} + (1 - q)\,\tl{k-1}{m}{p}$. Since
$\tl{k-1}{m-1}{p} \ge \tl{k-1}{m}{p}$, this expression increases with $q$, and at $q = p$ it
equals $\tl{k}{m}{p}$.
\end{proof}

\begin{theorem}[Soundness]
\label{thm:sound}
Every derivable judgment is valid.
\end{theorem}

\begin{proof}
By induction on the derivation. Fix $\theta$, $E$, and $h$, and write $\Pr$ for the
probability under the semantics of the expression in the conclusion.

\textsc{Ret} and \textsc{Fun}: the outcome is a value that satisfies the type, so both
probabilities are $0$.

\textsc{Oracle}: the outcome is a value $u$ drawn from $E(o, v, h)$. It is bad when
$\neg\psi_o(v, u)$, which has probability at most $\delta_o$ because $E \in \cO$.

\textsc{Let}: let $(r_1, h_1)$ be the outcome of $e_1$, and let $q$ be the probability that
$r_1$ is bad for $\varphi_1$. If the final outcome is bad, then either $r_1$ is bad, or $r_1$
is a value $v$ that satisfies $\varphi_1$ and the outcome of $e_2[v/x]$ from $h_1$ is bad.
For such a $v$, the substitution $\theta$ extended with $v$ for $x$ satisfies
$\Phi, \varphi_1$, so the induction hypothesis for $e_2$ bounds the conditional probability
of the second event by $\varepsilon_2$, for each $v$ and $h_1$. The value $r_1$ satisfies
$\varphi_1$ with probability at most $1 - q$. Hence
$\Pr[\Bad] \le q + (1 - q)\,\varepsilon_2$, which increases with $q$ and is at most
$\varepsilon_1 \oplus \varepsilon_2$ since $q \le \varepsilon_1$. The same argument with
failures in place of bad outcomes gives $\Pr[\Fail] \le \gamma_1 \oplus \gamma_2$, since the
final outcome is $\bot$ when $r_1 = \bot$.

\textsc{Check}: let $r$ be the outcome of $e$. The final outcome is bad only if $r$ is a
value that violates $\varphi$ and the checker accepts it. Such a value violates $p_j$ by the
first side condition and hence $\chi$ by the second, so $r$ is bad for $\chi$, which has
probability at most $\varepsilon$. A value that violates $p_j$ is accepted with probability
at most $\alpha_j$, whatever the history. So $\Pr[\Bad] \le \alpha_j \varepsilon$. The final
outcome is a failure if $r$ is a failure for $\chi$, or if $r$ satisfies $\chi$ and the
checker rejects it. Let $f \le \gamma$ be the probability of the first event. A value that
satisfies $\chi$ satisfies $p_j$ and is rejected with probability at most $\beta_j$. So
$\Pr[\Fail] \le f + (1 - f)\,\beta_j \le \gamma \oplus \beta_j$. In the first event we have
counted a value that violates $\chi$ and satisfies $\varphi$ as a failure although the checker
may accept it; this only loosens the bound.

\textsc{Check-Any}: a bad final outcome is a value that violates $p_j$ and was accepted.

\textsc{OrElse}: let $(r_1, h_1)$ be the outcome of $e_1$, let $q$ be the probability that
$r_1$ is bad, and let $a$ be the probability that $r_1 = \bot$. Then $q \le \varepsilon_1$
and $q + a \le \gamma_1$. The final outcome is bad when $r_1$ is bad, or $r_1 = \bot$ and the
outcome of $e_2$ from $h_1$ is bad. The induction hypothesis for $e_2$ holds for every $h_1$,
so
\[
\Pr[\Bad] \le q + a\,\varepsilon_2 \le q + (\gamma_1 - q)\,\varepsilon_2
  = q\,(1 - \varepsilon_2) + \gamma_1 \varepsilon_2
  \le \varepsilon_1 + (\gamma_1 - \varepsilon_1)\,\varepsilon_2 .
\]
In the same way $\Pr[\Fail] \le q + a\,\gamma_2 \le \varepsilon_1 + (\gamma_1 -
\varepsilon_1)\,\gamma_2$.

\textsc{Vote} and \textsc{Vote-Any}: let $X_i$ be $1$ if the outcome of the $i$-th run is bad.
The $i$-th run starts from the history that the earlier runs left, and the induction
hypothesis for $e$ holds for every history, so the conditional probability of $X_i = 1$ is at
most $\varepsilon$ whatever happened before. If the final outcome is bad, a value that
violates $\varphi$ occurs more than $k/2$ times, so at least $\lfloor k/2 \rfloor + 1$ runs
are bad, and Lemma~\ref{lem:dom} bounds the probability. For the failures in \textsc{Vote},
let $X_i$ be $1$ if the $i$-th run fails. By the side condition all runs that do not fail
return the same value. If fewer than $\lceil k/2 \rceil$ runs fail, that value occurs more
than $k/2$ times and is the final outcome. So a failure of the vote needs at least
$\lceil k/2 \rceil$ failed runs.

\textsc{If}: the test is closed after substitution, so one branch runs, and $\theta$
satisfies the context of that branch.

\textsc{Sub}: a value that violates $\varphi'$ violates $\varphi$, so the bad outcomes and
the failures for $\varphi'$ are among those for $\varphi$.
\end{proof}

The proof uses the bounds on the environment at two places, the rules \textsc{Oracle} and
\textsc{Check}, and it uses them for the history at hand. This is why the theorem holds for
adaptive environments. The rules \textsc{Let}, \textsc{OrElse} and \textsc{Vote} use that the
induction hypothesis holds for every history.

\subsection{The Number of Calls}

Retries and votes buy reliability with calls. The grades also bound how many. For an
expression $e$, define $c(e)$ by recursion: $c(\ret\;v) = c(f(v)) = 0$; $c(o(v)) = 1$;
$c(\letin{x}{e_1}{e_2}) = c(e_1) + c(e_2)$; $c(\chk{j}{e}) = c(e) + 1$ for a checker that is
not exact and $c(e)$ for one that is; $c(\vote{k}{e}) = k \cdot c(e)$; the larger branch for
a conditional; and
\[
c(e_1 \orelse e_2) = c(e_1) + \gamma_1 \cdot c(e_2),
\]
where $\gamma_1$ is a bound on the failures of $e_1$ from a typing judgment. In a
$\kw{let}$, the judgments used for $e_2$ must not assume that $x$ satisfies its type: a wrong
value of $x$ reaches $e_2$ as well, and it may make $e_2$ run longer.

\begin{proposition}[Calls]
\label{prop:calls}
For every environment in $\cO$ and every history, the expected number of calls of oracles and
of checkers that are not exact in an evaluation of $e$ is at most $c(e)$.
\end{proposition}

\begin{proof}
By induction on $e$. The second branch of a fallback runs when the first abstains, which is
a failure and has probability at most $\gamma_1$ by Theorem~\ref{thm:sound}. The expectation
of a sum is the sum of the expectations, and the bound for $e_2$ holds for every history and
every value of the variables.
\end{proof}

\subsection{Exact Values and Tightness}
\label{sec:tight}

How much do the grades lose? To answer, we need the true worst case.

\paragraph{The game.} Fix a closed program, a closed type, and an objective: the probability
of a bad outcome, the probability of a failure, or the expected number of calls. The
supremum of the objective over $\cO$ is the value of a game in which the environment moves
at every call. Because an environment may depend on the history in any way, the game can be
solved from the end. Let $K(u)$ be the value of the rest of the program after an oracle call
returns $u$. The environment may put probability at most $\delta_o$ on wrong answers, so the
value before the call is
\[
\max_{u \text{ right}} K(u) \;+\; \delta_o \cdot
\max\Bigl(0,\; \max_{u \text{ wrong}} K(u) - \max_{u \text{ right}} K(u)\Bigr) .
\]
The environment gives its favourite right answer, and with probability $\delta_o$ its
favourite wrong one, if that is better for it. A checker is treated in the same way: the
environment chooses the probability of acceptance at one end of the interval that the rates
allow. For a finite set of values this is a finite computation. The artifact performs it
with exact rational arithmetic (Section~\ref{sec:tests}).

\paragraph{Where the grades are exact.} For the patterns of Section~\ref{sec:consequences}
the grades equal the values of the game: a single call, a checked call, a loop of retries
with an exact or an imperfect checker, a vote over calls with unique right answers, and two
unchecked calls in a row whose errors show in the result. The artifact confirms the equality
on \exCases{} instances. Among \exPrograms{} random programs, the bound on bad outcomes
equals the value of the game for \exTightBad{} percent, the bound on failures for
\exTightFail{} percent, and the bound on calls for \exTightCost{} percent. With sums in place
of the operation $\oplus$, as in the first version of this paper, the first two numbers are
\exUnionBad{} and \exUnionFail{} percent.

\paragraph{Why two numbers are not enough.} The grades cannot be exact for all programs. The
rule \textsc{OrElse} assumes the worst of its first branch twice: that it is wrong as often
as possible, and that it abstains as often as the difference of its grades allows. An
environment may be unable to do both. Consider
\[
\bigl(\letin{x}{\chk{j}{o_1(v)}}{o_2(x)}\bigr) \orelse o_3(v)
\]
with an exact checker $j$ for the answers of $o_1$, and failure rates $1/4$, $1/3$ and $1/2$.
The first branch is wrong when $o_2$ fails, with probability up to $1/3$, and it fails with
probability up to $1/4 \oplus 1/3 = 1/2$. The rule gives the bound
$1/3 + (1/2 - 1/3) \cdot 1/2 = \exGapBound$ on bad outcomes. But the first branch abstains
only when $o_1$ fails, and then $o_2$ does not run. An environment that lets $o_1$ fail with
probability $d$ makes the program wrong with probability $(1 - d)/3 + d/2$, which is at most
\exGapExact{}, at $d = 1/4$. The value of the game is \exGapExact{}. To be exact, a system
would have to describe the first branch by the set of pairs of an error probability and an
abstention probability that environments can produce, and not by one bound on each. We keep
two numbers. They are what a programmer states and reads, and the measurement above says how
often they suffice.

\section{Consequences}
\label{sec:consequences}

\subsection{A Sound Check Removes Error, a Retry Buys Results}

\begin{corollary}
\label{cor:loop}
Let $o$ be an oracle with specification $\psi_o$ and failure rate $\delta$, and let $j$ be an
exact checker with $p_j(x) \models \varphi$ and $\psi_o(v, x) \models p_j(x)$. Then
\[
\vdash \retry{n}{(\chk{j}{o(v)})} : \rty{x}{\varphi} \at{0}{\delta^n},
\qquad c = 1 + \delta + \cdots + \delta^{n-1} \le \frac{1}{1 - \delta} .
\]
\end{corollary}

\begin{proof}
By \textsc{Oracle} and \textsc{Check}, $\chk{j}{o(v)}$ has the grade $(0, \delta)$ and one
call. By \textsc{OrElse} and induction on $n$, the grade of $n$ copies is
$(0,\; \delta \cdot \delta^{n-1})$, and the calls are $1 + \delta \cdot c_{n-1}$.
\end{proof}

The loop never returns a wrong value, for any $n$ and any $\delta$, including $\delta = 1$.
The failure rate of the oracle affects only how often the loop gives up and how many calls it
makes. To give up with probability at most $\eta$, it suffices to take
$n \ge \log \eta / \log \delta$. All three bounds are attained: an environment that fails
with probability $\delta$ at every attempt, with an answer that the checker rejects, reaches
them. The corollary is the precise sense of the title. A checker turns the errors of the
oracle into abstentions, and a retry turns abstentions into calls.

\subsection{Error Comes from Unchecked Calls}

Let $U(e)$ be the multiset of oracle calls in $e$ that are not inside a check by a sound
checker, where a call inside $\vote{k}{}$ counts $k$ times.

\begin{proposition}
\label{prop:unchecked}
If $\Phi \vdash e : \rty{x}{\varphi} \at{\varepsilon}{\gamma}$ is derived without the rules
\textsc{Sub} and \textsc{Check-Any}, then
\[
\varepsilon \;\le\; 1 - \prod_{o \in U(e)} (1 - \delta_o) \;\le\; \sum_{o \in U(e)} \delta_o .
\]
In particular, a program whose oracle calls are all inside sound checks has $\varepsilon = 0$.
\end{proposition}

\begin{proof}
By induction on the derivation. The middle expression is the $\oplus$ of the rates in
$U(e)$. \textsc{Oracle} contributes $\delta_o$. In \textsc{Check} a sound checker gives $0$,
and any checker gives at most the error of its argument. \textsc{Let} combines the bounds of its parts with
$\oplus$, and $U$ is the union. For \textsc{OrElse},
$\varepsilon_1 + (\gamma_1 - \varepsilon_1)\,\varepsilon_2 \le \varepsilon_1 \oplus
\varepsilon_2$ because $\gamma_1 \le 1$. For a vote, more than half of $k$ runs are bad only
if some run is bad, which has probability at most the $\oplus$ of $k$ copies of
$\varepsilon$. \textsc{If} gives the common bound of its branches.
\end{proof}

Ten unchecked calls in a row, each wrong one time in twenty, are wrong together with
probability up to \exChainTen{} percent. Consider a pipeline that turns an informal claim
into a checked theorem. An oracle $\nm{formalize}$ with failure rate $\delta_f$ writes a
formal statement $s$. An oracle $\nm{prove}$ with failure rate $\delta_p$ writes a proof, and
a kernel checks it:
\[
\letin{s}{\nm{formalize}(c)}{\;\letin{\pi}{\retry{n}{(\chk{\nm{kernel}_s}{\nm{prove}(s)})}}
  {\;\ret\;(s, \pi)}} .
\]
Its grade is $(\delta_f,\; \delta_f \oplus \delta_p^n)$ for the specification that $s$ states
$c$ and $\pi$ proves $s$. The kernel removes $\delta_p$ from the error, however poor the
prover is. Nothing removes $\delta_f$: no checker accepts exactly the statements that say
what was meant. The error of the pipeline is the error of its unchecked step, and that step
is the one a human has to read. For a machine-checked library, what there is to read is the
audit surface of the statement~\citep{tpp1}.

\subsection{An Imperfect Check Is Amplified by Retries}
\label{sec:imperfect}

\begin{corollary}
\label{cor:imperfect}
Let $o$ have failure rate $\delta$, and let $j$ be a checker with rates $\alpha$ and $\beta$,
with $p_j(x) \models \varphi$ and $\psi_o(v, x) \models p_j(x)$. Let
$r = \delta\,(1 - \alpha) + (1 - \delta)\,\beta$. Then
\[
\vdash \retry{n}{(\chk{j}{o(v)})} : \rty{x}{\varphi} \at{\varepsilon_n}{\varepsilon_n + r^n},
\qquad
\varepsilon_n = \delta\alpha\,\frac{1 - r^n}{1 - r}
\;\xrightarrow[n \to \infty]{}\; \frac{\delta\alpha}{\delta\alpha + (1 - \delta)(1 - \beta)} .
\]
\end{corollary}

\begin{proof}
One attempt has the grade $(\delta\alpha,\; \delta \oplus \beta)$, and
$(\delta \oplus \beta) - \delta\alpha = r$. By \textsc{OrElse},
$\varepsilon_n = \delta\alpha + r\,\varepsilon_{n-1}$ and
$\gamma_n = \delta\alpha + r\,\gamma_{n-1}$, which gives the two sums. The limit follows from
$1 - r = \delta\alpha + (1 - \delta)(1 - \beta)$.
\end{proof}

The limit is the probability that an answer is wrong given that the checker accepted it: the
loop that runs until acceptance returns the first accepted answer. It depends on the
generator as much as on the checker. For a good generator, with $\delta$ small, it is about
$\delta\alpha$, far below the checker's own rate $\alpha$. For a weak generator it is far
above it. With $\delta = 0.8$, $\alpha = 0.05$ and $\beta = 0$, one attempt is wrong with
probability \exPatchOne{} percent, five attempts with \exPatchFive{} percent, and the limit is
\exPatchLimit{} percent, after \exPatchCalls{} attempts on average. The bounds are attained
by an environment that fails independently at each attempt.

Two readings follow. First, the number of retries is a parameter of correctness when the
checker is imperfect, and only a parameter of cost when the checker is sound. A harness that
raises its retry limit to get more answers gets more wrong answers with a test suite and none
with a kernel. Second, generated tests as checkers~\citep{chen2022} and models as judges have
rates that are themselves uncertain, and the corollary says what to estimate: the rate of
false acceptance on the wrong answers that this generator produces, since the limit is linear
in $\alpha$ for small $\alpha$ and the factor is $\delta/(1 - \delta)$. A companion paper
treats the same check against an attacker who controls the generator: there the rate of false
acceptance of a model that judges is $1$, and only code can endorse~\citep{tpp7}.

\subsection{A Vote Keeps Its Guarantee, with Conditions}
\label{sec:votes}

\begin{corollary}
\label{cor:vote}
Let $o$ have failure rate $\delta$ and a specification with a unique right answer for each
input. Then
\[
\vdash \vote{k}{o(v)} : \rty{x}{\psi_o(v, x)}
  \at{\tl{k}{\lfloor k/2 \rfloor + 1}{\delta}}{\tl{k}{\lceil k/2 \rceil}{\delta}},
\qquad c = k .
\]
\end{corollary}

The bounds are those of $k$ independent answers, and they hold for answers that are not
independent, by Lemma~\ref{lem:dom}. They are attained by an environment that fails
independently and always with the same wrong value. For $\delta = 0.3$ the bound on wrong
results is \exVoteOne{} for one call, \exVoteThree{} for $k = 3$, \exVoteFive{} for $k = 5$,
\exVoteEleven{} for $k = 11$, and \exVoteTwentyFive{} for $k = 25$.

Three conditions separate a vote from a check. A vote reduces error and does not remove it,
and each halving of the error costs a constant number of further calls, where a check costs
one. A vote amplifies whichever side has the majority: for $\delta = 0.6$ the same sequence
is \exVoteBadOne, \exVoteBadThree, \exVoteBadFive, \exVoteBadEleven, \exVoteBadTwentyFive. And
a vote needs the right answer to be unique. If two answers are both right, an environment
that alternates between them is never wrong and never reaches a majority, and
\textsc{Vote-Any} gives the failure bound $1$. Free-form answers are of this kind, which is
why votes in practice are taken over a normalized final answer~\citep{wang2022}.

\subsection{Retries and Votes Need a Fresh Chance}
\label{sec:fresh}

The class $\cO$ bounds the failure of an oracle for every history. A weaker assumption bounds
only the probability that a given call fails. The rules \textsc{OrElse} and \textsc{Vote} are
unsound for it.

\begin{proposition}
\label{prop:repeat}
Let $E$ answer the first call to $o$ with a value that the exact checker $j$ rejects with
probability $\delta$ and with a right value otherwise, and let it repeat its first answer at
every later call. Then every call to $o$ fails with probability $\delta$,
$\retry{n}{(\chk{j}{o(v)})}$ abstains with probability $\delta$ for every $n \ge 1$, and
$\vote{k}{o(v)}$ is wrong with probability $\delta$ for every $k \ge 1$.
\end{proposition}

\begin{proof}
All calls return the first answer. With probability $\delta$ the checker rejects it $n$
times and the loop abstains, and all $k$ runs of the vote agree on it.
\end{proof}

A deterministic oracle behaves like this: a language model at temperature zero gives the
same answer to the same question. Such an oracle is not in $\cO$ unless $\delta_o$ is $0$ or
$1$ for the inputs concerned. A retry or a vote helps only when each attempt has a fresh
chance to succeed, by new randomness or by a changed input, for example one that includes
the checker's report. Adaptivity is allowed; repetition is the case that the class excludes.

\section{Testing the Claims}
\label{sec:tests}

The artifact is one Python script, \texttt{grades.py}, of about 420 lines without
dependencies. Values are the numbers $0$ to $3$. There are four oracles with functional and
relational specifications, three checkers with positive rates, and exact checkers for
arbitrary sets of values. Predicates are sets, so entailment is inclusion.

\paragraph{Grades and game.} One function computes the grades of Figure~\ref{fig:rules} and
the bound $c$ of Proposition~\ref{prop:calls} for a closed program. A second function computes
the three values of the game of Section~\ref{sec:tight} by backward induction, with exact
fractions.

\paragraph{Tests.} For \exPrograms{} random programs with at most seven calls, each grade is
at least the corresponding value of the game. This tests Theorem~\ref{thm:sound} and
Proposition~\ref{prop:calls} against the strongest environment for each program, which no
sampling of environments would find. The script counts how often a grade equals the value,
which gives the percentages of Section~\ref{sec:tight}. It then breaks the rules one at a
time, in \exRules{} ways that make a grade smaller: an oracle with half its rate, a maximum
in place of $\oplus$ for either grade, a check that ignores its second side condition for
either grade or its rate of false rejection, a fallback that forgets the error of its first
branch in either grade, a vote with a higher threshold or without the condition of
uniqueness, and two variants of the bound on calls. Each broken rule is refuted by a random program, the slowest after \exSlowest{}
programs. Finally the script checks the closed forms of Section~\ref{sec:consequences} on
\exCases{} programs, where grades and values coincide, and the example of
Section~\ref{sec:tight}, where they do not. The command \texttt{make -C artifact check} runs
everything in about fifteen seconds and compares the output, including the numbers quoted in
this paper, with the committed copy.

\section{Related Work}
\label{sec:related}

\paragraph{Reliability and union bounds.} \citet{carbin2013} verify quantitative reliability
for programs on unreliable hardware. Their judgments bound the probability that a value is
correct, and operations fail independently. We separate wrong values from abstentions and
allow adaptive failures. \citet{barthe2016} give a program logic whose proofs compose failure
probabilities with the union bound, which is the content of a weaker form of our rule
\textsc{Let}. Our rules \textsc{Check}, \textsc{OrElse} and \textsc{Vote} are not union
bounds: the first removes or scales error, the second multiplies, and the third is a
binomial tail. The grades of our judgments are an instance of effect systems with a monoid
of effects, whose semantics graded monads describe~\citep{katsumata2014}; here the effect of
a computation is a pair of probabilities and sequencing is $\oplus$ in each component.

\paragraph{Checking instead of trusting.} \citet{blum1995} design programs that check their
work, \citet{mcconnell2011} survey algorithms that return certificates, and proof-carrying
code checks the proofs of an untrusted producer~\citep{necula1997}. In these works the
checked result is correct outright. Our grades say what a checker does to a probabilistic
producer, what remains when some steps have no checker, and what an imperfect checker costs.

\paragraph{Abstention.} In classification, the option to reject trades errors for
abstentions~\citep{chow1970}, and selective classifiers choose when to
abstain~\citep{geifman2017}. Our rule \textsc{Check} is that trade made by a test on the
answer, and with a sound test the error drops to zero.

\paragraph{Generate, test, vote.} Programs around language models use all three constructs.
Self-consistency takes a majority over sampled answers~\citep{wang2022}. CodeT generates
tests together with code and selects the code by the tests that it passes~\citep{chen2022},
which is a check with a generated checker. In program repair, patches that pass the tests and are wrong are
known as overfitting patches~\citep{smith2015}; Corollary~\ref{cor:imperfect} is the
arithmetic of that phenomenon for a loop. These works measure. We give the bounds that hold
for every environment within the rates.

\paragraph{This series.} The same check, read for an attacker in place of an unreliable
model, turns an injection into a choice among accepted values~\citep{tpp7}. The calls that
retries add are paid for in tokens, and a companion paper bounds that
cost~\citep{tpp6}.

\section{Limits and Open Problems}
\label{sec:limits}

\paragraph{Limits.} The failure rates are assumptions about the world. The type system does
not justify them; it says what follows from them. The class $\cO$ asks for a fresh chance at
every call, which a deterministic oracle does not give. The bounds are worst cases over
adaptive environments and may be far from the behavior of a particular model. The side
conditions are semantic, and a checker for these types needs a logic for them, as refinement
type systems do. The language has no recursion. The grades are not exact in general
(Section~\ref{sec:tight}).

\paragraph{Open problems.}
\begin{enumerate}
\item \emph{Exact abstractions.} Describe a computation by the set of pairs of error and
abstention probabilities that environments can produce. Is this set always a polygon with
few vertices, and do the rules extend to it?
\item \emph{Rates with confidence.} A rate is estimated from samples. A judgment should then
hold with a stated confidence, and the estimate needed for
Corollary~\ref{cor:imperfect} is the rate of false acceptance on the generator's own wrong
answers, which is the hardest to sample.
\item \emph{Recursion.} With loops that run until success, the number of calls is a random
quantity, and the bound on calls becomes a fixed point.
\item \emph{One budget.} Calls cost tokens~\citep{tpp6}. A single judgment could give the
error, the abstention, and the price of a program, and tell which of a check, a retry, or a
vote buys the most reliability for a given budget.
\item \emph{Mechanization.} The proof of Theorem~\ref{thm:sound} is an induction over a
finite probabilistic semantics, and Lemma~\ref{lem:dom} is its only probabilistic argument.
\end{enumerate}

\section{Conclusion}

A sound checker turns the errors of an oracle into abstentions, and a retry turns abstentions
into calls. An imperfect checker lets a fraction of the errors through, and retries collect
that fraction again at every attempt. A vote lowers error without a checker when answers are
unique and usually right. The type system states this with two numbers, proves it for
oracles that adapt to the past, and is exact on the patterns that programs use. The error
that remains in a program is the error of the calls that no sound checker covers.

\bibliographystyle{ACM-Reference-Format}
\bibliography{refs}

\end{document}
